% DRAFT: replaces the equal-technology accounting as the baseline.
\paragraph{Informal technology.} Following \citet{IhrigMoe2004} and \citet{ElginOztunali2012}, all capital is formal: $Y_{f,t}=K_t^{\alpha}(X_tN_{f,t})^{1-\alpha}$ and $Y_{i,t}=A_i e^{z^r_t}X_tN_{i,t}^{\eta}$ with $0<\eta<1$. The justification is not that informal units use no capital, but what the household survey measures. The labor income reported by own-account workers and informal employers is mixed income or profit, and therefore already contains the return to the capital of their own unit; the wages they pay are counted in the income of their employees. Informal survey income is thus a measure of informal value added per worker. Formal survey wages, by contrast, measure only the labor share of formal value added, before the wedge. The informal output share is then
\begin{equation}
\theta_t=\frac{(1-\alpha)R_t\,s_t/(1-s_t)}{1+(1-\alpha)R_t\,s_t/(1-s_t)},
\label{eq:share-labor-only}
\end{equation}
where $R_t$ is informal value added per worker relative to the gross formal wage. Informal employees of firms with six or more workers are the exception: they work with the firm's capital, so their value added is their wage divided by $1-\alpha$; the earnings series is built accordingly (Section~\ref{sec:data}). Under equal technologies the factor $1-\alpha$ disappears and \eqref{eq:share-labor-only} becomes the share of labor income, which is therefore an upper bound.

\paragraph{Capital share.} The capital share of the formal sector is estimated rather than assumed. Capital is built by the perpetual-inventory method from gross fixed capital formation in constant prices, and $\ln(Y_f/N_f)=c+\alpha\ln(K/N_f)+gt+e$ is estimated by dynamic OLS \citep{StockWatson1993} on annual data for 1993--2025, with $Y_f$ computed from \eqref{eq:share-labor-only} so that $\alpha$ is a fixed point. {{ALPHA_RESULT}}
